% !Mode:: "TeX:UTF-8"

\begin{publication}
\noindent\textbf{（一）发表的相关论文}

\begin{publist}
\item 一作. IEEE Transactions on Affective Computing (\textbf{TAC}), 2026. （EI, SCI, IF: 11.3, 中科院1区TOP, JCR1区, CCF-B类期刊, 第5章）
\item 学生一作. International Journal of Computer Vision (\textbf{IJCV}), 2026, 134. （EI, SCI, IF: 10.3, 中科院2区，JCR1区, CCF-A类期刊, 第4章）
\item 一作. Expert Systems with Applications (\textbf{ESWA}), 2026, 324. （EI, SCI, IF: 9.4, 中科院1区TOP, JCR1区, CCF-B类期刊，第2章）
\item 一作. Neurocomputing, 2024, 575.（EI, SCI, IF: 6.7, 中科院2区, JCR1区, CCF-C类期刊, 第2章）
\item 一作. IEEE International Conference on Acoustics, Speech and Signal Processing (\textbf{ICASSP}), 2026. （EI, CCF-B类会议，第3章）
\end{publist}


\noindent\textbf{（二）在投的相关论文}

一作. Information Fusion, 2026. (EI, SCI, IF: 17.4, 中科院1区TOP, JCR1区). （第3章）

\noindent\textbf{（三）已获得的专利}

% ===== 以下为非盲审原始内容，保留备用 =====
学生一作, 发明专利, 2025.

\noindent\textbf{（四）参与的科研项目及获奖情况}

姚鸿勋等。农场环境自适应感知与认知技术研究，科技创新2030“新一代人工智能”重大项目。课题编号：2021ZD0110901。

% ===== 以下为非盲审原始内容，保留备用 =====
% \item 姚鸿勋等。农场环境自适应感知与认知技术研究，科技创新2030“新一代人工智能”重大项目。课题编号：2021ZD0110901。
%\vfill
%\hangafter=1\hangindent=2em\noindent
%\setlength{\parindent}{2em}
\end{publication}
